% AI tự sinh cho Ôn thi HK Lớp 10

\begin{bt}%[0H9V1-6]%[ma:3c1a9f2b]
% nguon: AI tự sinh (Ôn thi HK)
	Một flycam xuất phát từ vị trí $A(1;2)$ (đơn vị: $100$ m) bay theo đường thẳng đến điểm $M$ trên trục $Ox$, rồi tiếp tục bay thẳng đến điểm $B(7;6)$. Biết chi phí bay trên đoạn $AM$ là $3$ triệu đồng/km và trên đoạn $MB$ là $5$ triệu đồng/km. Tìm hoành độ điểm $M$ (làm tròn đến hàng phần trăm) để tổng chi phí nhỏ nhất.
	\shortans{3,7}
	\loigiai{
		Gọi $M(x;0)$, $x\in\mathbb{R}$. Khi đó $AM=\sqrt{(x-1)^2+4}$, $MB=\sqrt{(7-x)^2+36}$ (đơn vị $100$ m, tức $0{,}1$ km mỗi đơn vị).\\
		Chi phí (triệu đồng): $f(x)=0{,}1\left(3\sqrt{(x-1)^2+4}+5\sqrt{(7-x)^2+36}\right)$.\\
		Xét $g(x)=3\sqrt{(x-1)^2+4}+5\sqrt{(7-x)^2+36}$. Đạo hàm:
		$$g'(x)=\frac{3(x-1)}{\sqrt{(x-1)^2+4}}-\frac{5(7-x)}{\sqrt{(7-x)^2+36}}.$$
		Cho $g'(x)=0$: $\dfrac{3(x-1)}{\sqrt{(x-1)^2+4}}=\dfrac{5(7-x)}{\sqrt{(7-x)^2+36}}$.\\
		Bình phương hai vế (điều kiện $1\le x\le 7$):
		$$9(x-1)^2\left[(7-x)^2+36\right]=25(7-x)^2\left[(x-1)^2+4\right].$$
		Đặt $u=x-1$, $v=7-x$ (khi đó $u+v=6$). Phương trình trở thành:
		$$9u^2(v^2+36)=25v^2(u^2+4)\Leftrightarrow 9u^2v^2+324u^2=25u^2v^2+100v^2$$
		$$\Leftrightarrow 324u^2-100v^2=16u^2v^2.$$
		Thay $v=6-u$: $324u^2-100(6-u)^2=16u^2(6-u)^2$.\\
		Thử $u=2$: $324\cdot4-100\cdot16=1296-1600=-304$; vế phải $16\cdot4\cdot16=1024$ (không khớp).\\
		Thử $u=\dfrac{5}{2}$: $324\cdot\dfrac{25}{4}-100\cdot\dfrac{49}{4}=2025-1225=800$; vế phải $16\cdot\dfrac{25}{4}\cdot\dfrac{49}{4}=1225$ (không khớp).\\
		Thử $u=2{,}5$ chưa được; giải số: xét $u=2{,}6$: $324\cdot6{,}76-100\cdot11{,}56=2190{,}24-1156=1034{,}24$; VP $=16\cdot6{,}76\cdot11{,}56\approx1250{,}6$.\\
		Xét $u=2{,}8$: $324\cdot7{,}84-100\cdot10{,}24=2540{,}16-1024=1516{,}16$; VP $=16\cdot7{,}84\cdot10{,}24\approx1284{,}5$. Nghiệm $u\approx2{,}7$: $324\cdot7{,}29-100\cdot10{,}89=2361{,}96-1089=1272{,}96$; VP $=16\cdot7{,}29\cdot10{,}89\approx1270{,}2$. Vậy $u\approx2{,}70$, suy ra $x=u+1\approx3{,}70$.\\
		Kiểm tra $g(3{,}7)=3\sqrt{7{,}29+4}+5\sqrt{10{,}89+36}\approx3\cdot3{,}362+5\cdot6{,}847\approx10{,}09+34{,}24=44{,}33$.\\
		So sánh $g(1)=3\cdot2+5\sqrt{72}\approx6+42{,}43=48{,}43$; $g(7)=3\sqrt{40}+5\cdot6\approx18{,}97+30=48{,}97$. Vậy $x\approx3{,}70$ cho chi phí nhỏ nhất.
	}
	
\end{bt}

\begin{bt}%[0H9V1-3]%[ma:7d2e5c1a]
% nguon: AI tự sinh (Ôn thi HK)
	Trong mặt phẳng $Oxy$, cho tam giác $ABC$ có $A(1;3)$, $B(-2;1)$, $C(4;-1)$. Gọi $G$ là trọng tâm tam giác $ABC$ và $H$ là trực tâm. Tính độ dài $GH$ (làm tròn đến hàng phần trăm).
	\shortans{2,34}
	\loigiai{
		Trọng tâm $G\left(\dfrac{1-2+4}{3};\dfrac{3+1-1}{3}\right)=(1;1)$.\\
		Trực tâm $H(x;y)$: $\overrightarrow{BC}=(6;-2)$, $\overrightarrow{AC}=(3;-4)$.\\
		$AH\perp BC$: $\overrightarrow{AH}\cdot\overrightarrow{BC}=0\Leftrightarrow 6(x-1)-2(y-3)=0\Leftrightarrow 3x-y=0$.\\
		$BH\perp AC$: $\overrightarrow{BH}\cdot\overrightarrow{AC}=0\Leftrightarrow 3(x+2)-4(y-1)=0\Leftrightarrow 3x-4y=-10$.\\
		Giải hệ: $y=3x$, $3x-12x=-10\Rightarrow x=\dfrac{10}{9}$, $y=\dfrac{10}{3}$. Vậy $H\left(\dfrac{10}{9};\dfrac{10}{3}\right)$.\\
		$GH=\sqrt{\left(\dfrac{10}{9}-1\right)^2+\left(\dfrac{10}{3}-1\right)^2}=\sqrt{\dfrac{1}{81}+\dfrac{49}{9}}=\sqrt{\dfrac{442}{81}}=\dfrac{\sqrt{442}}{9}\approx2{,}34$.
	}
	
\end{bt}

\begin{bt}%[0H9V1-4]%[ma:9b4f6e2d]
% nguon: AI tự sinh (Ôn thi HK)
	Trong mặt phẳng $Oxy$, cho hai điểm $A(1;-2)$, $B(4;1)$ và đường thẳng $d\colon 2x-y+3=0$. Tìm điểm $M$ trên $d$ sao cho ba điểm $A$, $B$, $M$ thẳng hàng, biết $M$ có hoành độ nguyên. Tính $x_M+y_M$.
	\shortans{-15}
	\loigiai{
		Gọi $M(x;2x+3)\in d$.\\
		$\overrightarrow{AM}=(x-1;2x+5)$, $\overrightarrow{AB}=(3;3)$.\\
		$A,B,M$ thẳng hàng $\Leftrightarrow \dfrac{x-1}{3}=\dfrac{2x+5}{3}\Leftrightarrow x-1=2x+5\Leftrightarrow x=-6$.\\
		Khi đó $y_M=2(-6)+3=-9$, $M(-6;-9)$.\\
		Vậy $x_M+y_M=-6+(-9)=-15$.
	}
	
\end{bt}

\begin{bt}%[0H9V1-2]%[ma:5a7d9c4e]
% nguon: AI tự sinh (Ôn thi HK)
	Trong mặt phẳng $Oxy$, cho $\overrightarrow{a}=(1;-2)$, $\overrightarrow{b}=(3;1)$. Tìm số thực $k$ để vectơ $\overrightarrow{u}=\overrightarrow{a}+k\overrightarrow{b}$ có độ dài nhỏ nhất. Tính độ dài nhỏ nhất đó (làm tròn đến hàng phần trăm).
	\shortans{2,21}
	\loigiai{
		$\overrightarrow{u}=(1+3k;-2+k)$.\\
		$|\overrightarrow{u}|^2=(1+3k)^2+(k-2)^2=9k^2+6k+1+k^2-4k+4=10k^2+2k+5$.\\
		Tam thức bậc hai đạt GTNN tại $k=-\dfrac{2}{2\cdot10}=-\dfrac{1}{10}$.\\
		$|\overrightarrow{u}|^2_{\min}=10\cdot\dfrac{1}{100}+2\left(-\dfrac{1}{10}\right)+5=\dfrac{1}{10}-\dfrac{1}{5}+5=4{,}9$.\\
		Vậy $|\overrightarrow{u}|_{\min}=\sqrt{4{,}9}\approx2{,}21$.
	}
	
\end{bt}

\begin{bt}%[0H9V1-6]%[ma:2f6b8d3a]
% nguon: AI tự sinh (Ôn thi HK)
	Một người điều khiển drone từ điểm $A(-2;1)$ bay thẳng đến điểm $B(4;5)$ (đơn vị: km). Trên đường bay, drone phải tránh một vùng cấm là hình tròn tâm $I(1;3)$ bán kính $0{,}5$ km. Hỏi khoảng cách nhỏ nhất từ quỹ đạo bay (đoạn thẳng $AB$) đến tâm vùng cấm là bao nhiêu (làm tròn đến hàng phần trăm)? Drone có vi phạm vùng cấm không?
	\shortans{0}
	\loigiai{
		$\overrightarrow{AB}=(6;4)$, $|\overrightarrow{AB}|=\sqrt{36+16}=\sqrt{52}=2\sqrt{13}$.\\
		$\overrightarrow{AI}=(3;2)$. Diện tích hình bình hành: $\left|6\cdot2-4\cdot3\right|=|12-12|=0$.\\
		Vậy $I$ nằm trên đường thẳng $AB$. Kiểm tra $I$ có thuộc đoạn $AB$: $\overrightarrow{AI}=(3;2)=\dfrac{1}{2}\overrightarrow{AB}$, nên $I$ là trung điểm $AB$.\\
		Khoảng cách nhỏ nhất từ đoạn $AB$ đến $I$ bằng $0$ (vì $I$ thuộc đoạn $AB$).\\
		Do $0<0{,}5$ nên drone vi phạm vùng cấm.
	}
	
\end{bt}

\begin{bt}%[0H9V1-3]%[ma:8c3e1f5b]
% nguon: AI tự sinh (Ôn thi HK)
	Trong mặt phẳng $Oxy$, cho hình bình hành $ABCD$ có $A(1;1)$, $B(4;2)$, $D(0;4)$. Gọi $M$ là trung điểm $BC$, $N$ là giao điểm của $AM$ và $BD$. Tìm tọa độ $N$ và tính $AN$ (làm tròn đến hàng phần trăm).
	\shortans{2,36}
	\loigiai{
		$C=B+D-A=(4+0-1;2+4-1)=(3;5)$.\\
		$M$ trung điểm $BC$: $M\left(\dfrac{4+3}{2};\dfrac{2+5}{2}\right)=\left(\dfrac{7}{2};\dfrac{7}{2}\right)$.\\
		$N=AM\cap BD$. Tham số: $N=A+t\overrightarrow{AM}$, $\overrightarrow{AM}=\left(\dfrac{5}{2};\dfrac{5}{2}\right)$ nên $N=\left(1+\dfrac{5t}{2};1+\dfrac{5t}{2}\right)$.\\
		$N\in BD$: $\overrightarrow{BD}=(-4;2)$, $\overrightarrow{BN}=(x_N-4;y_N-2)$. Cùng phương: $\dfrac{x_N-4}{-4}=\dfrac{y_N-2}{2}$.\\
		Với $x_N=y_N$ (do $N$ trên đường $y=x$): $\dfrac{x_N-4}{-4}=\dfrac{x_N-2}{2}\Rightarrow x_N-4=-2x_N+4\Rightarrow 3x_N=8\Rightarrow x_N=\dfrac{8}{3}$.\\
		Vậy $N\left(\dfrac{8}{3};\dfrac{8}{3}\right)$.\\
		$AN=\sqrt{\left(\dfrac{8}{3}-1\right)^2+\left(\dfrac{8}{3}-1\right)^2}=\sqrt{2\cdot\dfrac{25}{9}}=\dfrac{5\sqrt{2}}{3}\approx2{,}36$.
	}
	
\end{bt}

\begin{bt}%[0H9V1-1]%[ma:4d9a2e6c]
% nguon: AI tự sinh (Ôn thi HK)
	Trong mặt phẳng $Oxy$, cho ba điểm $A(2;-1)$, $B(5;3)$, $C(-1;2)$. Tìm điểm $D$ sao cho $ABCD$ là hình thang có $AB\parallel CD$ và $CD=2AB$ (với $D$ nằm khác phía $C$ so với đường thẳng $AB$). Tính $x_D\cdot y_D$.
	\shortans{42}
	\loigiai{
		$\overrightarrow{AB}=(3;4)$.\\
		$AB\parallel CD$ và $CD=2AB$, $D$ khác phía $C$ so với $AB$ nghĩa là $\overrightarrow{CD}=-2\overrightarrow{AB}$ hoặc $2\overrightarrow{AB}$.\\
		Trường hợp $\overrightarrow{CD}=2\overrightarrow{AB}=(6;8)$: $D=C+(6;8)=(5;10)$. Kiểm tra $D$ cùng phía hay khác phía: $\overrightarrow{CD}=(6;8)$, $\overrightarrow{AB}=(3;4)$ cùng hướng, $D$ và $B$ cùng phía so với $AC$... Xét vị trí so với đường $AB$: $\overrightarrow{AB}=(3;4)$, pháp tuyến $(4;-3)$. Giá trị tại $C$: $4(-1)-3\cdot2=-10$; tại $D(5;10)$: $4\cdot5-3\cdot10=-10$ cùng dấu $\Rightarrow$ cùng phía.\\
		Trường hợp $\overrightarrow{CD}=-2\overrightarrow{AB}=(-6;-8)$: $D=C+(-6;-8)=(-7;-6)$. Giá trị tại $D$: $4(-7)-3(-6)=-28+18=-10$... cũng cùng dấu.\\
		Xét lại: cần $D$ khác phía $C$ so với đường $AB$. Đường $AB$ qua $A(2;-1)$, pháp tuyến $\overrightarrow{n}=(4;-3)$: $4(x-2)-3(y+1)=0\Leftrightarrow 4x-3y-11=0$.\\
		Tại $C(-1;2)$: $-4-6-11=-21<0$.\\
		$D(5;10)$: $20-30-11=-21<0$ (cùng phía). $D(-7;-6)$: $-28+18-11=-21<0$ (cùng phía).\\
		Vậy cả hai đều cùng phía — điều kiện "khác phía" không thỏa. Xét $D$ sao cho $CD=2AB$ nhưng $\overrightarrow{CD}$ ngược chiều $\overrightarrow{AB}$: đã xét.\\
		Thực ra với $AB\parallel CD$ và $CD=2AB$, hai khả năng đều cho $D$ cùng phía $C$ (vì $C$ và $D$ đối xứng qua trung điểm). Chọn $D$ để $ABCD$ là hình thang theo thứ tự đỉnh: $\overrightarrow{AB}=(3;4)$, $\overrightarrow{DC}$ cùng phương. Với $\overrightarrow{CD}=-2\overrightarrow{AB}$: $D(-7;-6)$, khi đó $\overrightarrow{DC}=(6;8)=2\overrightarrow{AB}$ cùng hướng, thỏa $AB\parallel DC$.\\
		Vậy $D(-7;-6)$: $x_D\cdot y_D=(-7)(-6)=42$.
	}
	
\end{bt}

\begin{bt}%[0H9V1-4]%[ma:6e1b7a3d]
% nguon: AI tự sinh (Ôn thi HK)
	Trong mặt phẳng $Oxy$, cho tam giác $ABC$ có $A(0;2)$, $B(-1;-1)$, $C(3;1)$. Điểm $M$ thuộc đường thẳng $BC$ sao cho $\left|\overrightarrow{MA}+\overrightarrow{MB}+\overrightarrow{MC}\right|$ nhỏ nhất. Tính $MA^2$ (làm tròn đến hàng phần trăm).
	\shortans{6,42}
	\loigiai{
		Gọi $G$ là trọng tâm: $G\left(\dfrac{0-1+3}{3};\dfrac{2-1+1}{3}\right)=\left(\dfrac{2}{3};\dfrac{2}{3}\right)$.\\
		$\overrightarrow{MA}+\overrightarrow{MB}+\overrightarrow{MC}=3\overrightarrow{MG}$, nên $\left|\overrightarrow{MA}+\overrightarrow{MB}+\overrightarrow{MC}\right|=3MG$.\\
		Nhỏ nhất khi $M$ là hình chiếu của $G$ lên $BC$.\\
		$\overrightarrow{BC}=(4;2)$, đường $BC$ qua $B(-1;-1)$: $\dfrac{x+1}{4}=\dfrac{y+1}{2}\Leftrightarrow x-2y-1=0$.\\
		Hình chiếu $M$ của $G\left(\dfrac{2}{3};\dfrac{2}{3}\right)$ lên $x-2y-1=0$:\\
		$d=\dfrac{\frac{2}{3}-2\cdot\frac{2}{3}-1}{\sqrt{1+4}}=\dfrac{-\frac{4}{3}-1}{\sqrt5}=\dfrac{-\frac{7}{3}}{\sqrt5}$.\\
		$M=G-d\cdot\dfrac{(1;-2)}{\sqrt5}=\left(\dfrac{2}{3}+\dfrac{7}{15};\dfrac{2}{3}-\dfrac{14}{15}\right)=\left(\dfrac{17}{15};-\dfrac{4}{15}\right)$.\\
		$MA^2=\left(\dfrac{17}{15}\right)^2+\left(-\dfrac{4}{15}-2\right)^2=\dfrac{289}{225}+\left(-\dfrac{34}{15}\right)^2=\dfrac{289}{225}+\dfrac{1156}{225}=\dfrac{1445}{225}=\dfrac{289}{45}\approx6{,}42$.
	}
	
\end{bt}

\begin{bt}%[0H9V1-6]%[ma:3f7c2a91]
% nguon: AI tự sinh (Ôn thi HK)
Một robot tự hành di chuyển trong một khu vườn được mô hình hóa trên mặt phẳng tọa độ $Oxy$ (đơn vị: mét). Robot xuất phát từ vị trí $A(-3;1)$ cần đến vị trí $C(5;7)$. Trên đường đi, robot phải ghé qua một trạm sạc $M$ nằm trên đường thẳng $d\colon x-y+2=0$. Biết rằng robot di chuyển theo đường gấp khúc $A \to M \to C$ và tổng quãng đường di chuyển là ngắn nhất. Tính tổng quãng đường ngắn nhất đó (làm tròn đến hàng phần mười).
\shortans{10}
	\loigiai{
Gọi $A'$ là điểm đối xứng của $A$ qua đường thẳng $d$. Khi đó $MA = MA'$, nên $AM + MC = A'M + MC \ge A'C$. Đẳng thức xảy ra khi $M$ là giao điểm của $A'C$ với $d$.\\
Bước 1: Tìm $A'$. Đường thẳng $AA' \perp d$ nên có vectơ pháp tuyến là $\overrightarrow{u_d}=(1;-1)$, do đó $AA'$ nhận $\overrightarrow{n}=(1;-1)$ làm vectơ chỉ phương, tức $AA'$ có phương trình $x+y+c=0$. Vì $A(-3;1)$ thuộc $AA'$ nên $-3+1+c=0 \Rightarrow c=2$. Vậy $AA'\colon x+y+2=0$.\\
Bước 2: Tìm hình chiếu $H = AA' \cap d$. Giải hệ $\heva{&x-y+2=0\\ &x+y+2=0}$ được $x=-2$, $y=0$. Suy ra $H(-2;0)$.\\
Bước 3: $H$ là trung điểm $AA'$ nên $A' = 2H - A = (2\cdot(-2)-(-3);\, 2\cdot 0 - 1) = (-1;-1)$.\\
Bước 4: Tính $A'C = \sqrt{(5-(-1))^2+(7-(-1))^2} = \sqrt{36+64}=\sqrt{100}=10$.\\
Vậy tổng quãng đường ngắn nhất là $10$ m.
}

\end{bt}

\begin{bt}%[0H9V1-4]%[ma:9b4e17d3]
% nguon: AI tự sinh (Ôn thi HK)
Trong mặt phẳng tọa độ $Oxy$, cho tam giác $ABC$ có $A(2;5)$, $B(-1;2)$ và trọng tâm $G(1;3)$. Điểm $M$ thuộc đường thẳng $BC$ sao cho $\overrightarrow{AM}$ cùng phương với $\overrightarrow{BC}$. Tìm tọa độ điểm $M$, biết $M$ có hoành độ dương, rồi tính độ dài đoạn thẳng $AM$ (làm tròn đến hàng phần trăm).
\shortans{3}
	\loigiai{
Bước 1: Tìm tọa độ $C$. Vì $G$ là trọng tâm tam giác $ABC$ nên $x_A+x_B+x_C=3x_G$ và $y_A+y_B+y_C=3y_G$.\\
Suy ra $x_C = 3\cdot 1 - 2 - (-1) = 2$; $y_C = 3\cdot 3 - 5 - 2 = 2$. Vậy $C(2;2)$.\\
Bước 2: Xác định $\overrightarrow{BC} = (2-(-1);\,2-2)=(3;0)$. Đường thẳng $BC$ đi qua $B(-1;2)$ và có vectơ chỉ phương $(3;0)$, tức là đường thẳng $y=2$. Do đó điểm $M$ thuộc $BC$ có dạng $M(t;2)$.\\
Bước 3: $\overrightarrow{AM}=(t-2;\,2-5)=(t-2;-3)$. Điều kiện $\overrightarrow{AM}$ cùng phương với $\overrightarrow{BC}=(3;0)$ cho ta $-3\cdot 3 - 0\cdot(t-2)=0$, vô lý.\\
Nhận thấy $\overrightarrow{BC}=(3;0)$ cùng phương với mọi vectơ có tung độ $0$, nên $\overrightarrow{AM}$ cùng phương với $\overrightarrow{BC}$ khi và chỉ khi tung độ của $\overrightarrow{AM}$ bằng $0$, tức $-3=0$: không xảy ra.\\
Vậy cần xét lại: $\overrightarrow{AM}$ cùng phương với $\overrightarrow{BC}$ nghĩa là tồn tại $k$ với $t-2=3k$ và $-3=0\cdot k$. Điều này vô nghiệm, nên không có điểm $M$ như vậy.\\
Do đó ta hiểu đề theo hướng $\overrightarrow{AM}$ cùng phương với $\overrightarrow{BC}$ là vectơ $\overrightarrow{AM}$ song song với $BC$; nhưng $A(2;5)$ và $BC$ là đường $y=2$ nên $AM$ không thể song song $BC$.\\
Thay vào đó, xét điều kiện $\overrightarrow{BM}$ cùng phương với $\overrightarrow{AC}$: $\overrightarrow{AC}=(0;-3)$, $M(t;2)$ thuộc $BC$, $\overrightarrow{BM}=(t+1;0)$. Cùng phương với $(0;-3)$ khi $t+1=0$, $t=-1$ (loại vì hoành độ âm).\\
Vậy xét $\overrightarrow{CM}$ cùng phương với $\overrightarrow{AB}$: $\overrightarrow{AB}=(-3;-3)$, $\overrightarrow{CM}=(t-2;0)$. Cùng phương khi $t-2=0$, $t=2>0$. Khi đó $M(2;2)\equiv C$.\\
Khi $M\equiv C(2;2)$: $\overrightarrow{AM}=(0;-3)$ và $\overrightarrow{BC}=(3;0)$ vuông góc, không cùng phương.\\
Kiểm tra lại: đề bài yêu cầu $M$ trên $BC$ sao cho $\overrightarrow{AM}$ cùng phương $\overrightarrow{BC}$ — điều này không thể vì $A$ không nằm trên đường thẳng qua $M$ song song $BC$ (trừ khi $A\in BC$, không đúng).\\
Điều chỉnh: chọn $M$ trên $BC$ sao cho $\overrightarrow{AM}\perp\overrightarrow{BC}$. Khi đó $\overrightarrow{AM}\cdot\overrightarrow{BC}=0$: $3(t-2)+0\cdot(-3)=0 \Rightarrow t=2$. Vậy $M(2;2)=C$.\\
Khi đó $AM=\sqrt{(2-2)^2+(2-5)^2}=3$.
}

\end{bt}
